\section*{第 \ref{ch:b3:canonical-ensemble} 章——正则系综}

\begin{solution}{exo:b3:canonical-ensemble:1}
(a) $\mathcal P_2/\mathcal P_1 = (g_2/g_1)\eu^{-\beta\epsilon}$。
(b) $\beta\epsilon = 2.1/(8.62\times10^{-5} \times 2500) = 9.7$：
占比 $3\eu^{-9.7} \approx \num{2e-4}$。(c) 火焰中
钠原子数约为 $\sim\num{e15}$ 个 / cm$^3$：即使其中仅 $10^{-4}$
的原子每隔数纳秒循环一次，即每秒可发出 $10^{19}$ 个黄色光子
——足以令人目眩。(d) $3\eu^{-\beta\epsilon} = 0.1$：
$T = \epsilon/(k_{\text{B}}\ln30) \approx \qty{7200}{K}$——恒星
光球层，而非火焰。
\end{solution}

\begin{solution}{exo:b3:canonical-ensemble:2}
(a) $z = 1 + \eu^{-\beta\epsilon} + \eu^{-2\beta\epsilon}$。(b)
$\langle E\rangle = \epsilon(\eu^{-\beta\epsilon} +
2\eu^{-2\beta\epsilon})/z$：$\to 0$（$T$ 低时），$\to \epsilon$
（平均能级）于高温时。(c) $\mathcal P_1 = 1/(\eu^{\beta
\epsilon} + 1 + \eu^{-\beta\epsilon})$ 随
$T$ 单调增大，趋近其上确界 $1/3$，仅在 $T \to \infty$ 时。(d)
玻尔兹曼权重在 $T > 0$ 时仅随能量减小；
粒子数反转需要\omterm{ex:b3:microcanonical-ensemble:negative}{负温度}（见
\cref{ex:b3:microcanonical-ensemble:negative}），任何
能量库都无法达到。
\end{solution}

\begin{solution}{exo:b3:canonical-ensemble:3}
(a) 对 $E_p = mgh$ 取玻尔兹曼因子，温度 $T$ 处处均匀。(b) $h_0 =
k_{\text{B}}T/mg = \qty{8.4}{km}$。(c) $\eu^{-8848/8440} \approx
0.35$：三分之一个大气压——这正是登顶者携带氧气的原因。
(d) 真实大气层并非等温：温度随
高度下降，因此分布偏离单一指数函数。
\end{solution}

\begin{solution}{exo:b3:canonical-ensemble:4}
(a) 三个平动\omterm{def:b3:lagrangian-mechanics:coordinates}{自由度}：$\tfrac32 R = \qty{12.5}{J/(mol.K)}$，
与实验精确吻合。(b) 加上两个转动\omterm{def:b3:lagrangian-mechanics:coordinates}{自由度}：$\tfrac52 R =
\qty{20.8}{J/(mol.K)}$，与测量相符：振动被冻结。(c)
$\tfrac72 R = \qty{29.1}{J/(mol.K)}$ 为预测值；实测值
$\approx26$ 表明振动仅部分解冻（$\theta_{\text{振动}}
\approx \qty{3400}{K}$）。(d) $3R = \qty{24.9}{J/(mol.K)}$：
杜隆–珀替定律，多数金属在室温下服从，
而金刚石则不然——正是冻结现象（见
\cref{exo:b3:canonical-ensemble:6}）。
\end{solution}

\begin{solution}{exo:b3:canonical-ensemble:5}
(a) $z = \eu^{-\beta\hbar\omega/2}/(1 - \eu^{-\beta\hbar\omega})$。
(b) $\langle E\rangle = -\partial_\beta\ln z$ 给出所述
形式，其中 $\langle n\rangle = 1/(\eu^{\beta\hbar\omega} - 1)$。(c)
$C \to k_{\text{B}}$（$T$ 高时）；$C \approx k_{\text{B}}(\beta
\hbar\omega)^2\eu^{-\beta\hbar\omega} \to 0$（低温时）。(d) 若
$\hbar\omega$ 为光量子能量，则 $\langle n\rangle$
为每个模式中的热光子数：普朗克定律
仅隔一章。
\end{solution}

\begin{solution}{exo:b3:canonical-ensemble:6}
(a) 数值上，$C = \tfrac12 \times 3Nk_{\text{B}}$ 出现在
$T/\theta_{\text{E}} \approx 0.34$ 附近。(b) $\theta_{\text{E}}/T =
4.3$：$C/3Nk_{\text{B}} = 4.3^2\eu^{-4.3}/(1 - \eu^{-4.3})^2
\approx 0.26$——金刚石在室温下的热容仅约为
经典热容的四分之一：所谓“反常”实为
显而易见的量子冻结。(c) $\theta_{\text{E}} = \qty{90}{K}$
使铅在 \qty{300}{K} 时深处于经典区。(d) 问题出在
单一频率假设：实际固体具有一系列模式，
一直延伸到长波声波，其低能量子给出 $T^3$
定律——德拜的修补，在\cref{ch:b3:photons-phonons}中。
\end{solution}

\begin{solution}{exo:b3:canonical-ensemble:7}
(a) $f(v) \propto v^2\eu^{-mv^2/2k_{\text{B}}T}$：$v_{\text{p}} =
\sqrt{2k_{\text{B}}T/m}$：对 N$_2$ 为 \qty{413}{m/s}，
对 H$_2$ 为 \qty{1540}{m/s}。(b) $mv_{\text{逃逸}}^2/2k_{\text{B}}T
\approx 730$（N$_2$），$52$（H$_2$）。(c) $\eu^{-730}$
绝不可能；每个停留时间内 $\eu^{-52} \sim \num{e-23}$，逸散很慢——
但炽热的上层大气（$\sim\qty{1000}{K}$）使
氢的指数降至 $\sim15$：氢气在
地质时间内逸散，氮气则留存。(d) 逃逸速率与外逸层
温度：木星的 \qty{60}{km/s} 势阱使氢的
指数也大得惊人——气态巨行星留住了小而温暖的行星
所会失去的气体。
\end{solution}

\begin{solution}{exo:b3:canonical-ensemble:8}
(a) $\partial_\beta^2\ln Z = \langle E^2\rangle - \langle
E\rangle^2$，且 $C = \partial_T\langle E\rangle =
-k_{\text{B}}\beta^2\partial_\beta\langle E\rangle$。(b) 两边
均给出 $Nk_{\text{B}}(\beta\epsilon)^2\eu^{\beta\epsilon}/(
\eu^{\beta\epsilon} + 1)^2$。(c) $\langle E\rangle \propto N$，
$\Delta E \propto \sqrt N$。(d) 系统对加热的响应强弱
等于其能量自发涨落的大小——
响应与涨落是同一个二阶
导数的两种读法，这一模式（涨落–耗散）
贯穿整个物理学。
\end{solution}

\begin{solution}{exo:b3:canonical-ensemble:9}
(a) $\ln2 = E_a\,\Delta T/k_{\text{B}}T^2$，其中 $\Delta T =
\qty{10}{K}$，$T = \qty{298}{K}$：$E_a =
0.693\,k_{\text{B}}T^2/10 \approx \qty{0.55}{eV}$。(b) 从
$298$ 降至 \qty{278}{K}：因子 $\eu^{E_a(1/278 - 1/298)/k_{\text{B}}}
\approx 4.4$ 倍慢——这正是冰箱保存食物的原因。(c)
$373 \to \qty{366}{K}$：速率降低 $\approx1.4$ 倍：
十一分钟的鸡蛋需要一刻钟。(d) 反应
由越过势垒的分子的指数尾部决定：稍稍改变温度，
尾部分子数便
极大地变化——平均值几乎无足轻重。
\end{solution}

\begin{solution}{exo:b3:canonical-ensemble:10}
(a) $f_{\text{折叠}} = 1/(1 + \eu^{-\beta(\Delta E -
T\Delta S)})$，其中展开态的熵已并入其
\omterm{prop:b3:canonical-ensemble:machine}{自由能}。(b) 当 $T_{\text{m}} = \Delta E/\Delta S$ 时，两
\omterm{prop:b3:canonical-ensemble:machine}{自由能}相等：各占一半。(c) $T_{\text{m}} =
\num{4.8e-19}/\num{1.38e-21} = \qty{348}{K}$（\qty{75}{\celsius}）；
宽度 $\delta T \sim k_{\text{B}}T_{\text{m}}^2/\Delta E \approx
\qty{3.5}{K}$。(d) 协同性使 $\Delta E$ 与
$\Delta S$ 均乘以同时断开的接触数，保持
$T_{\text{m}}$ 不变，而宽度 $\propto 1/\Delta E$ 缩小：DNA
链在几开尔文范围内分离，这正是
PCR 仪能在“熔解”与“退火”之间干净循环的原因。
\end{solution}

\begin{solution}{exo:b3:canonical-ensemble:11}
(a) $z = 2\cosh(\beta\mu B)$; $M = N\mu\tanh(\beta\mu B) \approx
N\mu^2B/k_{\text{B}}T$（宗量较小时）：居里定律 $1/T$。(b)
$\mu_{\text{B}}B/k_{\text{B}}T$：\qty{300}{K} 时为 $\num{2.2e-3}$；
在 \qty{1}{K} 时为 $0.67$——由无所偏向变为强烈取向。(c)
$S$ 仅取决于 $\mu B/k_{\text{B}}T$；在熵固定时降低 $B$，
$T$ 随之成比例下降：$\qty{1}{K} \to
\qty{0.1}{K}$。(d) 当 $k_{\text{B}}T$ 达到自旋–自旋
能量时，熵不再由磁场控制：
$\mu_0\mu_{\text{B}}^2/4\pi a^3 \approx \qty{7e-26}{J} \sim
\qty{5}{mK}$——经典下限（核磁矩弱约一千倍，
将其推至微开尔文量级）。
\end{solution}

\begin{solution}{exo:b3:canonical-ensemble:12}
(a) $z_{\text{转动}} = \sum_J(2J+1)\eu^{-\beta BJ(J+1)}
\to \int(2J+1)\eu^{-\beta BJ(J+1)}\dd J = k_{\text{B}}T/B$：于是
$\langle E\rangle = k_{\text{B}}T$，从而 $C_{\text{转动}} =
k_{\text{B}}$。(b) H$_2$：$\theta_{\text{转动}} = \qty{87}{K}$——
台阶位于实验室范围内；N$_2$：\qty{2.9}{K}，
仅在液氦附近才冻结，因此氮气始终呈现
$\tfrac52 R$。(c) 平台 $\tfrac32 R$（低于 $\sim\qty{50}{K}$），
$\tfrac52 R$（从 $\sim200$ 到 $\sim\qty{1000}{K}$），上升
趋向 $\tfrac72 R$，靠近 $\theta_{\text{振动}} \approx
\qty{6000}{K}$ 附近。(d) 奇数与偶数的 $J$ 分属不同的核自旋
种类，二者转换缓慢，因此冷氢的 $C_V$
取决于其正–仲历史——核统计经由
量热计检验。
\end{solution}

\begin{solution}{pb:b3:canonical-ensemble:1}
\textbf{1.} $V = \tfrac43\pi r^3 = \qty{4.0e-20}{m^3}$: $m' = V
\Delta\rho = \qty{8.3e-18}{kg}$，重力 $m'g = \qty{8.1e-17}{N}$。
\textbf{2.} $n(h) = n_0\,\eu^{-m'gh/k_{\text{B}}T}$：这是
气压定律，缩小到了显微镜载玻片上。
\textbf{3.} $h_0 = k_{\text{B}}T/m'g = \num{4.04e-21}/\num{8.1e-17}
\approx \qty{50}{\micro m}$。
\textbf{4.} $h_0 \propto 1/r^3$：半径 $10\%$ 的散布，意味着
标高 $30\%$ 的散布——多分散颗粒把
阶梯抹成一团糊。佩兰借助
反复离心分级了数月。
\textbf{5.} 同一公式，质量相差 $10^{10}$ 倍：颗粒的
“大气”变浅了 $10^{10}$ 倍——千米量级缩小
到数十微米，这正是使其能在显微镜下整体
可观察的原因。
\textbf{6.} 相邻比值 $0.55$、$0.55$、$0.57$：为常数，
在计数误差范围内——指数规律。$h_0 =
\qty{30}{\micro m}/\ln(100/55) \approx \qty{50}{\micro m}$。
\textbf{7.} $k_{\text{B}} = m'g\,h_0/T = \num{8.1e-17} \times
\num{5.0e-5}/293 \approx \qty{1.4e-23}{J/K}$。
\textbf{8.} $N_{\text{A}} = R/k_{\text{B}} \approx
\num{6.0e23}\,\unit{mol^{-1}}$。
\textbf{9.} 此处在百分之几以内（采用理想化数据）；
佩兰的真实实验数据围绕现代值散布 $\pm15\%$，
——人工计数颗粒而言已令人惊叹，而对
数量级的确定则是决定性的。
\textbf{10.} $m_{\text{H}} = \qty{e-3}{kg/mol}/N_{\text{A}} =
\qty{1.7e-27}{kg}$。
\textbf{11.} 宏观化学提供 $R$；光学显微观测
与称量提供 $m'$；计数提供阶梯：三项
台式测量交叉印证，共同确定了谁也看不见的
原子的质量。
\textbf{12.} 无：推导只用到了“系统与能量库交换能量”
这一条件——玻尔兹曼因子与尺寸无关，
这正是佩兰所验证的。
\textbf{13.} $D = k_{\text{B}}T/6\pi\eta r =
\num{4.04e-21}/\num{4.0e-9} \approx \qty{e-12}{m^2/s}$。
\textbf{14.} $\sqrt{2Dt} = \sqrt{\num{1.2e-10}} \approx
\qty{11}{\micro m}$ 每分钟：很容易
借助目镜测微尺与耐心测得。
\textbf{15.} 两种独立的现象，两个独立的公式，
同一个数：静态阶梯与动态抖动相互印证，
不给巧合留下余地——分子假说
预言了这两者。
\textbf{16.} $\sqrt{k_{\text{B}}T/m} = \sqrt{\num{4.04e-21}/
\num{4.8e-17}} \approx \qty{9}{mm/s}$ 每个分量。
\textbf{17.} 颗粒每秒被撞击 $10^{19}$ 次，
在纳米距离内即忘掉其速度：弹道飞行
不可观测，眼睛所见的是其积分——即
扩散性随机游走。
\textbf{18.} 斯托克斯阻力（本科二年级分册中流体的
黏度）提供 $6\pi\eta r$；\omterm{thm:b3:canonical-ensemble:boltzmann}{正则系综}提供
$k_{\text{B}}T$：爱因斯坦的 $D$ 即二者之商，力学与统计
融于同一个分式之中。
\textbf{19.} 虚构之物无法计数：一旦 $N_{\text{A}}$ 成为
两个带误差的测量数之比，原子便成为
实验对象——奥斯特瓦尔德于 1909 年公开承认。
\textbf{20.} 天蓝色的散射、用测得的电子电荷做的电解、
镭中积累的氦：四条互不相关的
物理途径汇聚于同一个 $6\times10^{23}$，使得该数
成为自然的性质，而非任何理论的产物。
\textbf{21.} 引力势能把颗粒往下拉，
构型熵把颗粒往上扩散，二者以 $T$ 为兑换率相互交换：
阶梯\emph{正是} $F = E - TS$ 的最小值。
\textbf{22.} $h_0 \propto 1/g$：在 $10^5g$ 时，
$h_0 = k_{\text{B}}T/(m'\times10^5g) = \num{4.04e-21}/
\num{9.8e-17} \approx \qty{40}{\micro m}$，对于蛋白质——
即沉降平衡，佩兰的实验如今
已是生物化学系中的日常实验。
\textbf{23.} 从“过高”（$h_0 \gg$ 视场深度，看不出
梯度）到“过薄”（$h_0 \lesssim r$）：
可用半径大致在 $0.1$--$\qty{0.5}{\micro m}$ 之间——佩兰的选择并非
侥幸，而是精心设计。
\textbf{24.} 既然 $k_{\text{B}}$ 已由定义精确给定，
同一实验可用来标定颗粒（其尺寸或密度），或
检验装置：物理图像——玻尔兹曼阶梯——并未
改变，变的只是哪一个量作为未知量。
\textbf{25.} 对于 \qty{0.2}{\micro m} 的颗粒，其数量每
$\sim\qty{35}{\micro m}$ 减半；由 $n_0\eu^{-m'gh/k_{\text{B}}
T}$ 读出的阶梯给出 $N_{\text{A}} \approx \num{6e23}$：
阿伏伽德罗常数是逐颗计数得来的——而原子
之争在显微镜载物台上尘埃落定。
\end{solution}
